\section{Experiments}
This section presents the experimental results.
We first compare \method with existing baselines, then evaluate 
its flexibility across storage backends, and finally assess 
the contribution of each component and key design choices.

\subsection{Experimental Setup}
\label{ssec:experimental_setup}
\noindent\textbf{Datasets.} 
We evaluate \method on LoCoMo~\cite{maharana2024evaluating}, a benchmark for long-horizon conversational memory.
Following prior work~\cite{yan2025memory,fang2025lightmem}, we exclude the adversarial subset and focus on the reasoning-intensive QA setting.
More dataset details are provided in Appendix~\ref{appendix:dataset_details}.

\noindent\textbf{Baselines.} 
We compare against two passive baselines: \emph{Full Text} and 
\emph{Naive RAG}~\cite{gao2023retrieval}, and four active memory systems: 
\emph{LangMem}~\cite{LangMemBlog}, 
\emph{Mem0}~\cite{chhikara2025mem0}, 
\emph{A-Mem}~\cite{xu2025mem}, and 
\emph{LightMem}~\cite{fang2025lightmem}.
Additional baseline details are in Appendix~\ref{appendix:baseline_details}.

\noindent\textbf{Evaluation Protocol.} Following prior work~\cite{yan2025memory,chhikara2025mem0}, we report three metrics: token-level F1 (F1), BLEU-1 (B1), and LLM-as-a-Judge accuracy (ACC).
F1 and B1 measure lexical overlap with the reference answer; 
ACC measures semantic correctness via a judge model.
GPT-4o-mini~\cite{gpt4ocard} and Claude-Haiku-4.5~\cite{anthropic2025claudehaiku45} are used as the backbones for the Memory Manager, Meta-Thinker, and Query Reasoner.
To isolate memory construction quality from answer-generation capacity, we fix GPT-4o-mini as both the Answer Agent and the LLM judge across all experiments.
The retrieval budget is top-$30$ entries, the iterative refinement budget is $H{=}3$, and we generate $J{=}5$ probe QA pairs per session for self-evolution.
Additional implementation details are in Appendix~\ref{appendix:implementation_details}.

% \noindent\textbf{Evaluation Protocol.} Following prior work~\cite{yan2025memory,chhikara2025mem0}, we report three metrics: token-level F1 (F1), BLEU-1 (B1), and LLM-as-a-Judge accuracy (ACC).
% F1 and B1 measure lexical overlap with the reference answer; 
% ACC measures semantic correctness via a judge model.
% To ensure fair comparison, we fix the retrieval budget to top-$30$ entries and keep the answer-generation setting consistent across all methods.

% \noindent\textbf{Implementation Details.}
% GPT-4o-mini~\cite{gpt4ocard} and Claude-Haiku-4.5~\cite{anthropic2025claudehaiku45} 
% are used as the default backbone for the Memory Manager, 
% Meta-Thinker, and Query Reasoner.
% The iterative query refinement budget is $H{=}3$.
% To isolate memory construction quality from answer-generation 
% capacity, we fix GPT-4o-mini as both the Answer Agent and the 
% LLM judge across all experiments.
% For in-situ self-evolution, we generate $J{=}5$ probe QA pairs 
% per session using Claude-Opus-4.5~\cite{anthropic2025claudeopus45}, 
% retrieve top-$30$ entries for verification.
% All retrieval uses text-embedding-3-small~\cite{openai2024embedding3}.

\begin{table*}[t]
\centering
\small
\caption{Results on LoCoMo across four question categories (multi-hop, temporal, open-domain, single-hop). We report F1, B1, and ACC (\%). Best results are in \textbf{bold}. GPT-4o-mini and Claude-Haiku-4.5 are backbones; GPT-4o-mini is the answer agent. $\method_{\mathrm{LM}}$ uses LightMem as storage backend.}
\label{tab:locomo_category_results}
\resizebox{\textwidth}{!}{
\begin{tabular}{c|l|ccc|ccc|ccc|ccc|ccc}
\toprule
\multirow{2}{*}{\textbf{Model}} & \multirow{2}{*}{\textbf{Method}}
& \multicolumn{3}{c|}{\textbf{Multi-Hop}}
& \multicolumn{3}{c|}{\textbf{Temporal}}
& \multicolumn{3}{c|}{\textbf{Open-Domain}}
& \multicolumn{3}{c|}{\textbf{Single-Hop}}
& \multicolumn{3}{c}{\textbf{Overall}} \\
\cmidrule(lr){3-5} \cmidrule(lr){6-8} \cmidrule(lr){9-11} \cmidrule(lr){12-14} \cmidrule(lr){15-17}
& 
& \textbf{F1} & \textbf{B1} & \textbf{ACC}
& \textbf{F1} & \textbf{B1} & \textbf{ACC}
& \textbf{F1} & \textbf{B1} & \textbf{ACC}
& \textbf{F1} & \textbf{B1} & \textbf{ACC}
& \textbf{F1} & \textbf{B1} & \textbf{ACC} \\
\midrule

\multirow{6}{*}{\textbf{\rotatebox{90}{GPT}}}
& Full Text              & 29.41 & 21.16 & 43.75 & 29.95 & 19.33 & 51.35 & 18.25 & 19.56 & 61.54 & 41.45 & 29.96 & 74.29 & 34.13 & 24.63 & 61.18 \\
& Naive RAG              & 15.84 & 9.50 & 31.25 & 17.30 & 12.36 & 35.14 & 17.40 & 16.65 & 46.15 & 39.32 & 30.35 & 58.57 & 27.14 & 20.41 & 46.05 \\
& LangMem                & 12.55 & 9.22 & 25.00 & 15.23 & 11.53 & 21.62 & 14.91 & 14.03 & 38.46 & 23.52 & 17.59 & 35.71 & 18.46 & 14.05 & 30.26 \\
& A-Mem                  & 15.56 & 10.88 & 31.25 & 55.01 & 42.40 & 51.35 & 18.18 & 15.27 & 53.85 & 42.72 & 32.43 & 62.86 & 37.90 & 28.85 & 52.63 \\
& LightMem               & 33.74 & 29.33 & 65.62 & \textbf{59.76} & \textbf{51.12} & 78.38 & \textbf{31.85} & \textbf{24.23} & \textbf{76.92} & 43.88 & 34.68 & 78.57 & 44.58 & 36.66 & 75.66 \\
& $\method_{\mathrm{LM}}$ & \textbf{48.15} & \textbf{39.67} & \textbf{78.12} & 57.21 & 41.94 & \textbf{83.78} & 24.58 & 22.44 & \textbf{76.92} & \textbf{50.45} & \textbf{38.66} & \textbf{82.86} & \textbf{49.40} & \textbf{38.28} & \textbf{81.58} \\
\midrule

\multirow{6}{*}{\textbf{\rotatebox{90}{Claude-Haiku}}}
& Full Text              & 29.41 & 21.16 & 43.75 & 29.95 & 19.33 & 51.35 & 18.25 & 19.56 & 61.54 & 41.45 & 29.96 & 74.29 & 34.13 & 24.63 & 61.18 \\
& Naive RAG              & 15.84 & 9.50 & 31.25 & 17.30 & 12.36 & 35.14 & 17.40 & 16.65 & 46.15 & 39.32 & 30.35 & 58.57 & 27.14 & 20.41 & 46.05 \\
& LangMem                & 20.05 & 14.85 & 34.38 & 34.72 & 26.33 & 37.84 & 20.01 & 20.85 & 69.23 & 22.65 & 16.19 & 48.57 & 24.81 & 18.78 & 44.74 \\
& A-Mem                  & 15.79 & 10.32 & 28.13 & 56.41 & 43.23 & 54.05 & 16.34 & 17.76 & 38.46 & 38.37 & 27.98 & 65.71 & 36.12 & 27.10 & 52.63 \\
& LightMem               & {35.11} & {31.85} & 59.38 & 58.42 & \textbf{49.85} & \textbf{89.19} & \textbf{32.60} & 24.43 & 69.23 & 44.06 & \textbf{36.56} & 71.43 & 44.69 & \textbf{37.77} & 73.03 \\
& $\method_{\mathrm{LM}}$ & \textbf{35.38} & \textbf{32.48} & \textbf{65.62} & \textbf{59.25} & 44.66 & 83.78 & 28.59 & \textbf{26.86} & \textbf{84.62} & \textbf{45.31} & 35.85 & \textbf{77.14} & \textbf{45.10} & 36.53 & \textbf{76.97} \\
\bottomrule
\end{tabular}}
% \vspace{-0.5}
\end{table*}

\subsection{Main Comparison with Baselines}
\label{ssec:evaluation_main_results}
To evaluate \method, we compare it with baselines.
We use LightMem as the storage backend of \method, denoted by $\method_{\mathrm{LM}}$. GPT-4o-mini and Claude-Haiku-4.5 are the backbones. 
Other settings follow these in Sec.~\ref{ssec:experimental_setup}.

% To evaluate the effectiveness of \method, we compare \method with baselines. We use LightMem as the storage backend of \method, denoted by $\method_{\mathrm{LM}}$.
% % LightMem is the storage backend of \method and 
% GPT-4o-mini and Claude-Haiku-4.5 are the backbones. 
% Other settings are the same as these in Sec.~\ref{ssec:experimental_setup}.

Table~\ref{tab:locomo_category_results} reports the results. Three findings emerge:
\textbf{(i)}~\emph{$\method_{\mathrm{LM}}$ achieves the best overall performance under both backbones}.
Under GPT-4o-mini, it reaches $49.40$ F1, $38.28$ B1, and $81.58$ ACC, improving over LightMem by $+4.82$ F1, $+1.62$ B1, and $+5.92$ ACC.
Under Claude-Haiku-4.5, it again achieves the best overall ACC, improving from $73.03$ to $76.97$ over LightMem.
\textbf{(ii)}~\emph{The gains are strong at the category level.}
Under GPT-4o-mini, $\method_{\mathrm{LM}}$ improves most on Multi-Hop and Single-Hop, raising ACC from $65.62$ to $78.12$ and from $78.57$ to $82.86$, respectively.
The Multi-Hop gains are consistent with diagnosis-guided iterative retrieval helping recover distributed evidence, while the Single-Hop gains suggest that construction guidance and self-evolution help preserve precise answer-bearing details.
% \textbf{(ii)}~\emph{The gains are also strong at the category level}. 
% Under GPT-4o-mini, $\method_{\mathrm{LM}}$ improves most on Multi-Hop and Single-Hop, raising ACC from $65.62$ to $78.12$ and from $78.57$ to $82.86$, respectively.
% Under Claude-Haiku-4.5, the largest improvement appears on Open-Domain, where ACC increases from $69.23$ to $84.62$.
% These patterns are consistent with the full design of $\method_{\mathrm{LM}}$, which combines reasoning-aware coordination over construction and retrieval with in-situ self-evolving memory repair. 
\textbf{(iii)}~\emph{$\method_{\mathrm{LM}}$ improves an already strong baseline.}
LightMem is already the strongest baseline, yet $\method_{\mathrm{LM}}$ further improves it under both backbones, suggesting that the gain comes from memory-cycle coordination rather than a stronger storage backend.
% We further analyze its flexibility across storage backends in Sec.~\ref{ssec:flexibility_storage}.

\subsection{Flexibility across Storage Backends}
\label{ssec:flexibility_storage}
To assess the flexibility of \method across storage backends, we instantiate it on top of three memory systems: Single-Agent~\cite{yan2025memory} ($\method_{\mathrm{SA}}$), 
A-Mem ($\method_{\mathrm{AM}}$), and 
LightMem ($\method_{\mathrm{LM}}$).
All other components and settings are fixed as in Sec.~\ref{ssec:experimental_setup}.

% To assess the flexibility of \method across storage backends, we instantiate it on top of three memory systems: a single-agent memory pipeline~\cite{yan2025memory} ($\method_{\mathrm{SA}}$), A-Mem ($\method_{\mathrm{AM}}$), and LightMem (($\method_{\mathrm{SA}}$)).
% All other components and settings are fixed as in Sec.~\ref{ssec:experimental_setup}.
% Across all settings, the Meta-Thinker, Query Reasoner, and Answer Agent are kept fixed, and only the storage backend used by the Memory Manager is changed.
% All other settings are the same as those in Sec.~\ref{ssec:experimental_setup}.

Table~\ref{tab:backend_results} reports results on LoCoMo under GPT-4o-mini.
Two observations emerge.
\textbf{(i)}~\emph{\method consistently improves all backends.}
In terms of ACC, \method improves the Single-Agent backend from $52.60$ to $84.87$, A-Mem from $52.63$ to $78.29$, and LightMem from $75.66$ to $81.58$.
For A-Mem and LightMem, the gains are also consistent in F1 and B1.
For the weaker Single-Agent backend, B1 decreases even though Acc rises sharply, suggesting that \method improves semantic correctness more than lexical overlap in this setting.
These results indicate that \method improves long-horizon memory across diverse storage implementations.
\textbf{(ii)}~\emph{The gains of \method complement storage quality rather than replace it.}
Among the enhanced variants, $\method_{\mathrm{LM}}$ achieves the strongest overall performance, which is consistent with LightMem being the strongest standalone backend.
This pattern suggests that \method improves how memory is coordinated, rather than relying on a particular storage design.

% Table~\ref{tab:backend_results} reports the results on LoCoMo with GPT-4o-mini as the backbone model. We observe: 
% \textbf{(i)}~\emph{\method consistently improves all backends.}
% For example, \method improves A-Mem from 52.63 to 78.29 ACC and 
% LightMem from 75.66 to 81.58 ACC, with gains spanning F1, B1, 
% and ACC across both backends.
% This confirms that \method's gains come from coordination over 
% the memory cycle, not from a specific storage implementation.
% \textbf{(ii)} \emph{The gains are complementary to storage quality}. Among the enhanced variants, $\method_{\mathrm{LM}}$ achieves the strongest overall performance, which is consistent with LightMem being the strongest standalone backend.
% This pattern indicates that the gains of \method complement storage quality rather than replace it.
% Stronger backends remain stronger after enhancement, while still benefiting from the proposed coordination policy.

\begin{table}[t]
\centering
\small
\caption{Flexibility across backends on LoCoMo under GPT-4o-mini. Best results per backend are in \textbf{bold}.}
\label{tab:backend_results}
{
\begin{tabular}{lccc}
\toprule
\textbf{Method} & \textbf{F1} & \textbf{B1} & \textbf{ACC} \\
\midrule
Single-Agent  & 22.64 & \textbf{17.24} & 52.60 \\
$\method_{\mathrm{SA}}$ & \textbf{23.64} & 12.94 & \textbf{84.87} \\
\midrule
A-Mem & 37.90 & 28.85 & 52.63 \\
$\method_{\mathrm{AM}}$ & \textbf{46.23} & \textbf{35.13} & \textbf{78.29} \\
\midrule
LightMem & 44.58 & 36.66 & 75.66 \\
$\method_{\mathrm{LM}}$ & \textbf{49.40} & \textbf{38.28} & \textbf{81.58} \\
\bottomrule
\end{tabular}}
\vspace{-0.5em}
\end{table}


\subsection{In-depth Dissection of MemMA}
\label{ssec:ablation}
\noindent\textbf{Ablation Studies.}
To understand the contributions of key components in \method, we implement three ablated variants on the Single-Agent backend: 
\textbf{(i)} \method/C removes Meta-Thinker guidance during construction and directly uses the Memory Manager for memory writing;
\textbf{(ii)} \method/R removes iterative retrieval, reverting to one-shot retrieval based on semantic similarity; and
\textbf{(iii)} \method/E removes the probe-and-repair loop of in-situ self-evolving memory construction and directly commits $M_{\tau}^{(0)}$ to the memory bank.
% Single-Agent is used as the storage backend. 

% GPT-4o-mini and Claude-Haiku-4.5 are used as the construction backends. 
% All other settings are the same as those in Sec.~\ref{ssec:experimental_setup}. 

Fig.~\ref{fig:ablation} reports the results under GPT-4o-mini and Claude-Haiku-4.5. The full $\method_{\mathrm{SA}}$ achieves the strongest overall performance, while the variants reveal complementary weaknesses. Specifically:
(i)~\emph{Iterative retrieval is the most critical forward-path component.}
$\method_{\mathrm{SA}}$/R causes the largest drop under both backbones, reducing ACC from $84.87$ to $70.39$ with GPT-4o-mini and from $88.82$ to $81.58$ with Claude-Haiku-4.5.
This confirms that one-shot retrieval remains a major bottleneck and that diagnosis-guided refinement is essential for narrowing the information gap.
(ii)~\emph{Self-evolution repairs construction omissions.}
$\method_{\mathrm{SA}}$/E causes the second-largest degradation 
% (ACC: $84.87$ $\to$ $73.68$ with GPT-4o-mini, $88.82$ $\to$ $84.21$ 
% with Claude-Haiku-4.5).
(ACC: $84.87$ $\to$ $73.68$ with GPT-4o-mini).
The large ACC drop with only moderate F1 change suggests that 
self-evolution mainly improves semantic correctness by repairing 
missing information during construction.
(iii)~\emph{Construction guidance reduces upstream noise.}
$\method_{\mathrm{SA}}$/C reduces ACC from  $88.82$ to $83.55$ with 
Claude-Haiku-4.5.
This shows that construction decisions benefit from explicit strategic guidance rather than local heuristics alone, as the Meta-Thinker helps determine what should be retained, consolidated, or resolved before information enters the memory bank.
These ablations confirm that \method's gains come from 
complementary improvements on both paths of the memory cycle.

\begin{figure}[t]
    \small
    \centering
    \begin{subfigure}{0.235\textwidth}
        \includegraphics[width=0.98\linewidth]{figures/single_ablation_gpt4o_mini.pdf}
        \vskip -0.5em
        \caption{GPT-4o-mini}
    \end{subfigure}
    \begin{subfigure}{0.235\textwidth}
        \includegraphics[width=0.98\linewidth]{figures/single_ablation_haiku45.pdf}
        \vskip -0.5em
        \caption{Claude-Haiku-4.5}
    \end{subfigure}
    \vskip -0.5em
    \caption{Ablation studies of $\method_{\mathrm{SA}}$ under GPT-4o-mini and Claude-Haiku-4.5 on LoCoMo.}
    \vskip -1em
    \label{fig:ablation}
\end{figure}

\begin{figure}[t]
    \small
    \centering
    \begin{subfigure}{0.235\textwidth}
        \includegraphics[width=0.98\linewidth]{figures/impact_retrieval_budget_LM.pdf}
        \vskip -0.5em
        \caption{$\method_{\mathrm{LM}}$}
    \end{subfigure}
    \begin{subfigure}{0.235\textwidth}
        \includegraphics[width=0.98\linewidth]{figures/impact_retrieval_budget_SA.pdf}
        \vskip -0.5em
        \caption{$\method_{\mathrm{SA}}$}
    \end{subfigure}
    \vskip -0.5em
    \caption{Impact of retrieval budget $k$ of \method under both GPT-4o-mini and Claude-Haiku-4.5.}
    \vskip -1em
    \label{fig:impact_retrieval_budget}
\end{figure}

\begin{figure}[t]
    \small
    \centering
    \begin{subfigure}{0.235\textwidth}
        \includegraphics[width=0.98\linewidth]{figures/impact_refinement_budget_gpt4omini.pdf}
        \vskip -0.5em
        \caption{GPT-4o-mini}
    \end{subfigure}
    \begin{subfigure}{0.235\textwidth}
        \includegraphics[width=0.98\linewidth]{figures/impact_refinement_budget_haiku45.pdf}
        \vskip -0.5em
        \caption{Claude-Haiku-4.5}
    \end{subfigure}
    \vskip -0.5em
    \caption{Impact of refinement budget $H$ of \method.}
    \vskip -1em
    \label{fig:impact_retrieval_refinement}
\end{figure}

\noindent\textbf{Impact of retrieval budget $k$.} 
We vary $k \in \{10,20,30,40,50\}$ on both Single-Agent and LightMem backends and report results in Fig.~\ref{fig:impact_retrieval_budget}. We observe that the optimal $k$ depends on storage quality.
% We vary $k \in \{10,20,30,40,50\}$ with GPT-4o-mini as the backbone. Fig.~\ref{fig:impact_retrieval_budget} reports results on both Single-Agent and LightMem backends. We observe that the optimal $k$ depends on storage quality.
For $\method_{\mathrm{LM}}$, ACC peaks at $k{=}30$--$40$ ($81.58$) and declines at $k{=}50$ ($79.61$), indicating a sweet spot beyond which additional retrieval introduces noise.
For $\method_{\mathrm{SA}}$, ACC increases steadily from $75.66$ at $k{=}10$ to $84.21$ at $k{=}50$, without saturation.
We attribute this contrast to storage quality: stronger backends produce higher-quality, less redundant entries, so a moderate $k$ suffices and excess retrieval dilutes the evidence; weaker backends need a larger $k$ to retrieve enough evidence from sparser memory banks.

\noindent\textbf{Impact of retrieval refinement budget $H$.}
We vary the refinement budget $H \in \{0,1,2,3,4,5\}$ under both GPT-4o-mini and Claude-Haiku-4.5.
The results of $\method_{\mathrm{SA}}$ and $\method_{\mathrm{LM}}$ are reported in Fig.~\ref{fig:impact_retrieval_refinement}.
We observe that ACC improves sharply from one-shot retrieval ($H{=}0$) to a small $H$ and then declines.
For example, $\method_{\mathrm{SA}}$'s ACC rises from $78.95$ at $H{=}0$ to $85.53$ at $H{=}2$, then drops back to $81.58$ at $H{=}4$.
This shows that diagnosis-guided refinement converges quickly: one or two additional retrieval rounds suffice to close most information gaps, while further iterations risk retrieval drift. This validates the effectiveness of the Meta-Thinker's answerability diagnosis, which directs each refinement step toward the specific missing evidence rather than redundant searches. More analysis of the impact of probe generation model are in Appendix~\ref{appendix:impact_probe_generation_model}.

% \subsection{Case Studies}
% \label{ssec:case_studies}
% To better understand why \method improves long-horizon QA, we examine representative cases from both paths of the memory cycle.
% Our findings indicate that:
% \textbf{(i)} on the forward path, Meta-Thinker guidance helps preserve important answer-bearing details during memory construction, and diagnosis-guided iterative retrieval recovers specific facts (e.g., entity names, temporal anchors) that initial top-$k$ retrieval misses;
% \textbf{(ii)} on the backward path, 
% in-situ self-evolution converts local probe failures into targeted memory repairs that transfer to downstream QA, for example by inserting missing named entities, sharpening vague event descriptions, and completing partial evidence clusters.
% Detailed examples and traces are in Appendix~\ref{appendix:case_studies_details}.

\subsection{Case Studies}
\label{ssec:case_studies}
We conduct a case study to better understand why \method improves long-horizon QA.
% To better understand why \method improves long-horizon QA, we examine representative cases from both paths of the memory cycle.
Our findings indicate that:
\textbf{(i)} on the forward path, construction-time Meta-Thinker guidance determines whether answer-bearing details survive in memory, while diagnosis-guided iterative retrieval determines whether missing evidence is surfaced before the system commits to an answer.
Importantly, iterative retrieval cannot compensate for details that were never preserved during construction.
The cases also show that the retrieval controller and the storage backend play distinct roles: the Meta-Thinker and Query Reasoner identify the information gap, while the backend determines whether the required evidence can actually be recovered;
\textbf{(ii)} on the backward path, in-situ self-evolution converts local probe failures into targeted memory repairs that transfer to downstream QA, for example by inserting missing named entities, sharpening vague event descriptions, and completing partial evidence clusters.
Detailed examples are in Appendix~\ref{appendix:case_studies_details}.